\cotefile{31}
\scanpage{33}

% B01 | upper right, heading and parenthesis
Calculs divers

(acyclicité \unsure{pointée|ponctuée|finitude} etc)

% B02 | upper left of centre, faint count
4 p Deligne ? $+ 3 + 6$

% archive sheet number: 32

\scanpage{35}

% B01 | top centre
\mspage{\unsure{5|S}}

% B02 | upper right
GRÖTHENDIECK.

% B03 | printed letterhead, upper left
% letterhead
\begin{center}
UNIVERSITÉ DE PARIS\\
---$\cdot$---\\
ÉCOLE NORMALE SUPÉRIEURE\\
---$\cdot$---\\
45, RUE D'ULM\\
ODÉON 06$\cdot$45\\
---$\cdot$---
\end{center}

% B04 | theorem, left label and seven lines
\uuline{Th 1} Soit $E$ un ensemble simplicial, \deleted{\illegible{1}} $G$ un groupe discret qui opère sur $E$ \deleted{\illegible{1}}par des transformations \unsure{simplicials|simpliciales}, $B\inserted{=E/G}$ l'ensemble simplicial quotient, \deleted{\illegible{1}} $p:E\to B$ la projection canonique. Si $E$ est connexe, l'image de $p_*:\pi_1(E)\longrightarrow\pi_1(B)$ est un sousgroupe distingué de $\pi_1(B)$, et $\pi_1(B)/\operatorname{im}p_*(\pi_1(E))=G/S$, où $S$ est le sousgroupe (distingué) de $G$ engendré par les stabilisateurs des sommets de $E$, et $E/S$ est le revêtement \deleted{\unsure{universel|galoisien}} de $B$ correspondant à \inserted{l'image de $p_*(\pi_1(E))$}

% B05 | Déf., two lines under the theorem
\uuline{Déf.} une arête de $E$ est un 1-simplexe de $E$ muni d'un signe $+$ ou $-$. Si $\sigma$ est un 1-simplexe, on pose $\partial_1(+\sigma)=\partial_1(\sigma)$, $\partial_0(+\sigma)=\partial_0(\sigma)$, $\partial_1(-\sigma)=\partial_0(\sigma)$, $\partial_{\unsure{0|d}}(-\sigma)=\partial_1(\sigma)$.

% B06 | lower paragraph, Un chemin through suivants
Un chemin est une suite \inserted{d'arêtes} $\gamma=(\gamma_1,\ldots,\gamma_n)$ telle que $\partial_1(\gamma_{i+1})=\partial_0(\gamma_i)$. On pose $\partial_1(\gamma)=\partial_1(\gamma_1)$ et $\partial_0(\gamma)=\partial_0(\gamma_n)$. Un \deleted{point} \inserted{sommet} de base $x_0$ étant fixé dans $E$, un lacet est un chemin $\gamma$ tel que $\partial_1(\gamma)=\partial_0(\gamma)=x_0$. 2 chemins $\gamma$ et $\gamma'$ sont homotopes si leurs réalisations géométriques sont homotopes modulo leurs extrémités. Cette relation d'équivalence est compatible avec la loi de composition « mise bout à bout » des chemins, et est engendrée par les suivants :

% archive sheet number: 34

\scanpage{36}

% B01 | printed letterhead, upper left
% letterhead
\begin{center}
UNIVERSITÉ DE PARIS\\
---$\cdot$---\\
ÉCOLE NORMALE SUPÉRIEURE\\
---$\cdot$---\\
45, RUE D'ULM\\
ODÉON 06$\cdot$45\\
---$\cdot$---
\end{center}

% B02 | first three lines, crossed out as a block
\deleted{Soit $E$ \deleted{\unsure{un ensemble simplicial}} un C-w-complexe, $G$ un groupe discret, opérant sur $E$ par des transformations qui préservent la décomposition cellulaire, $B=E/G$ \deleted{l'ensemble} le C-w-complexe quotient de $E$}

% B03 | Lemme 1, two lines
\uline{Lemme 1} : $\forall x$ sommet de $E$, $\beta$ chemin de $B$ tel que $\partial_1(\beta)=p(x)$, $\exists\gamma$ chemin de $E$ tel que $\partial_1(\gamma)=x$ et $p(\gamma)=\beta$

% B04 | Dém, three lines
\uuline{Dém} Il suffit de le voir si $\beta$ est une arête. Soit $\lambda$ une arête de $B$ telle que $p\lambda=\beta$. On a \deleted{$\unsure{\partial}$} $p(\partial_1\lambda)=\partial_1(p\lambda)=\partial_1\beta=p(x)$, d'où $\exists g\deleted{\unsure{g}}\in G$ $x=g.\partial_1\lambda$, et $\gamma=g\lambda$ répond à la question

% B05 | Lemme 2, below centre fold
\uline{Lemme 2} : l'image de $p_*:\pi_1(E)\to\pi_1(B)$ est un sous groupe distingué de $\pi_1(B)$

% B06 | Dém, four lines
\uuline{Dém} : Soit $\gamma$ un lacet de \deleted{$\unsure{B}$} $E$, $\beta$ un lacet de $B$, $\lambda$ un chemin de $E$ tel que $\partial_1(\lambda)=x_0$ et $p_{\deleted{\unsure{*}}}(\lambda)=\beta$. On a $p(\partial_0\lambda)=\partial_0\beta=b_0=px_0$, d'où $\partial_0\lambda=gx_0$, et le lacet $(\lambda,g\gamma,-\lambda)$ se projette en $(\beta,p(\gamma),-\beta)$ ce qui démontre le lemme 2.

% B07 | underlined heading in lower quarter
\uline{Définition de $\chi:G\longrightarrow\pi_1(B)/\operatorname{Im}\pi_1(E)$ quand $E$ connexe.}

% B08 | final paragraph, five lines
Pour tout $g\in G$, soit $\gamma$ un chemin de $E$ tel que $\partial_1(\gamma)=x_0$, $\partial_0(\gamma)=gx_0$. \deleted{\unsure{soit $\chi(g)$ la classe de|et $\chi(g)$ la classe de}} $p(\gamma)$ est un lacet de $B$, soit $\chi(g)$ sa classe dans $\pi_1(B)/\operatorname{im}\pi_1(E)$. Elle ne dépend pas du chemin $\gamma$ choisi, car, si $\gamma'$ est un autre chemin, $(\gamma,-\gamma')$ est un lacet de $B$, et \deleted{\unsure{$p(\gamma)$ est homotope|par définition}} $p(\gamma)$ et $p(\gamma')$ ont même classe mod l'image de $\pi_1(E)$.

% archive sheet number: 35

\scanpage{37}

% B01 | printed letterhead, upper left
% letterhead
\begin{center}
UNIVERSITÉ DE PARIS\\
---$\cdot$---\\
ÉCOLE NORMALE SUPÉRIEURE\\
---$\cdot$---\\
45, RUE D'ULM\\
ODÉON 06$\cdot$45\\
---$\cdot$---
\end{center}

% B02 | first line
Il résulte du lemme 1 que $\chi$ est surjective.

% B03 | first dash, homomorphism paragraph
--- Montrons que $\chi$ est \deleted{\unsure{surj.|aussi}} un homomorphisme : \inserted{Soient $g_1,g_2\in G$} Soit $\gamma_1$ un \deleted{la} chemin tel que $\partial_1(\gamma_1)=x_0$, $\partial_0(\gamma_1)=g_1x_0$, idem \unsure{p\textsuperscript{r}|pr} $\gamma_2$ tel que $\partial_1(\gamma_2)=x_0$, $\partial_0(\gamma_2)=g_2x_0$. On a $\partial_1(g_1\gamma_2)=g_1x_0$ et $\partial_0(g_1\gamma_2)=g_1g_2x_0$ : le chemin $\inserted{\gamma=}(\gamma_1,g_1\gamma_2)$ vérifie $\partial_1\gamma=x_0$, $\partial_0\gamma=g_1g_2x_0$. \deleted{\unsure{Or on a|On a}} On a $p\gamma=(p\gamma_1,pg_1\gamma_2)=(p\gamma_1,p\gamma_2)$ \deleted{\unsure{or|donc}} d'où $\chi(g_1g_2)=\chi(g_1)\chi(g_2)$

% B04 | second dash and fixed-vertex construction
--- Le noyau de $\chi$ contient évidemment le stabilisateur de $x_0$. Soit $x$ un sommet de $E$ et $s\in G$ tel que $sx=x$. Soit $\deleted{\unsure{\gamma}}\lambda$ un chemin de $E$ tel que $\partial_1(\deleted{\gamma}\lambda)=x_0$, $\partial_0\deleted{\gamma}\lambda=x$. \deleted{$\unsure{(\gamma,-s\gamma)|(\lambda,-s\lambda)}$} $\gamma=(\lambda,-s\lambda)$ est un chemin de $E$ tel que $\partial_1(\gamma)=x_0$, $\partial_0\gamma=sx_0$. d'où $\chi(s)=\operatorname{classe}(p(\lambda,-s\lambda))=\operatorname{classe}(p(\lambda,-\lambda))=e\,(\text{élément neutre})$.

% B05 | paragraph beginning Le noyau, lower half
Le noyau de $\chi$ contient donc le sousgroupe $S$ \deleted{\unsure{des}} engendré par les stabilisateurs des \deleted{\unsure{points}} sommets de $E$. Ce sous groupe est distingué puisque le conjugué du stabilisateur de $x$ est le stabilisateur d'un transformé de $x$.

% B06 | third dash
--- On sait que $\operatorname{Ker}\chi\supset S$. Montrons qu'il ne le contient pas strictement.

% B07 | Lemme 3 at foot
\uline{Lemme 3} : $E/S$ est un revêtement \inserted{normal} de $E/G$, de groupe $G/S$.

% archive sheet number: 36

\scanpage{38}

% B01 | printed letterhead, upper left
% letterhead
\begin{center}
UNIVERSITÉ DE PARIS\\
---$\cdot$---\\
ÉCOLE NORMALE SUPÉRIEURE\\
---$\cdot$---\\
45, RUE D'ULM\\
ODÉON 06$\cdot$45\\
---$\cdot$---
\end{center}

% B02 | Dém., first four manuscript lines
\uuline{Dém}. $G/S$ opère sans point fixe sur $E/S$ : si $gx=s_0x$, $g=s_0s$ avec $sx=x\Rightarrow g\in S$. Si $g\notin S$, $g$ transforme tout sommet de $E$ en un sommet distinct, donc tout simplexe en un simplexe disjoint. $E/G={}^{E/S}\!/\!{}_{G/S}$.

% B03 | short rule and underlined conclusion
--- \uline{Le Théorème 1} \deleted{\unsure{I|1}} résulte de ce lemme et de ce qui a été vu précédemment.

% B04 | short line below conclusion
\deleted{$\unsure{\text{que }\pi_1(B)|p_*\pi_1(B)|p_*\pi_1(E)}$}

% archive sheet number: 37

\scanpage{39}

% B01 | printed letterhead, upper left
% letterhead
\begin{center}
UNIVERSITÉ DE PARIS\\
---$\cdot$---\\
ÉCOLE NORMALE SUPÉRIEURE\\
---$\cdot$---\\
45, RUE D'ULM\\
ODÉON 06$\cdot$45\\
---$\cdot$---
\end{center}

% B02 | cancelled line beside the address
\deleted{\unsure{En particulier si $E$ est}}

% B03 | Th 2, upper paragraph
\uuline{Th 2} : Si $G$ est un groupe fini d'ordre $n$, l'image de $p_*:H_*(E)\longrightarrow H_*(B)$ contient tous les éléments \deleted{d'ordre $n$} divisibles par $n$ dans $H_*(B)$.

% B04 | Dém, first two lines
\uuline{Dém} : Soit $\zeta$ \deleted{\unsure{un cycle}} $=\sum_{i\in I}a_i\sigma_i$ un cycle de $B$, où $\sigma_i$ sont des simplexes. Pour tout $\sigma_i$, soit $\tau_i$ simplexe de $E$ tel que $p\tau_i=\sigma_i$. Alors

% B05 | large sum and conclusion above fold
\[
\theta=\sum_{g\in G,\,i\in I}ga_i\tau_i\ \text{est un cycle de } E,\ \text{et }p(\theta)=\deleted{\unsure{\zeta|g}}n\zeta\ \text{\deleted{\unsure{est stable|est un}}}.
\]

% B06 | below right of sum
D'où le Th 2.

% B07 | Th 3, across and below fold
\uuline{Th 3} : Si \deleted{\unsure{$G$ est un|lorsque}} $E$ est \deleted{\unsure{acycle|acyclique}} contractile, $S={}$ ss-groupe engendré par les stabilisateurs fini d'ordre $n$, $\pi_i(B)$ est un groupe \deleted{\unsure{fini}} dont \inserted{tous les éléments sont} d'ordre qui divise une puissance de $n$ pour $i>1$

% B08 | Dém, last three lines
\uuline{Dém} : \deleted{\unsure{Comme|D'après}} D'après le Th. 1, $E/S$ est le revêtement universel de $B$, d'où $\pi_i(B)=\pi_i(E/S)$. D'après le Th. 2, \deleted{\unsure{par}} les $H_i(E/S)$

% archive sheet number: 38

\scanpage{42}
\mspage{\unsure{1}}

% B01 | head, lines 1–4
Soit $B$ un anneau avec unité et un \unsure{gpe|groupe} fini $G$ d'automorphismes $A=B^G$. On pose $\operatorname{Tr}_G b=\sum_{g\in G}gb$. Conditions équivalentes, ($q$ est un entier donné) : $A$ est un sous-anneau, et $\operatorname{Tr}_G:B\to A$ est un \unsure{hom} de bi-$A$-modules

% B02 | left of opening list, read clockwise
% margin note written sideways
\margin{\unsure{Voy. dans le Commentaire|Voy. dans le} \unsure{précis ou corollaire Voy. de Serre|précis au corollaire Voy. de Serre} Comment voir ça}

% B03 | four conditions at upper left
\[
\begin{aligned}
\text{a)}\quad & A=\operatorname{Tr}_G B\qquad\text{i.e. }\widehat H^0(G,B)=\{0\}\\
\text{a$'$)}\quad & 1=\operatorname{Tr}_G b\qquad\text{avec }b\in B\\
\text{b)}\quad & \widehat H(G,M)=0\qquad\text{pour }\unsure{\forall|\text{tout}}\ (B,G)\text{-module }M\\
\unsure{b(q)|b_q}\text{)}\quad & \widehat H^q(G,M)=0
\end{aligned}
\]

% B04 | Démonstration, first four lines
Démonstration : $a)\Longrightarrow a')$ trivial, $a'\Longrightarrow b$ car on a, en désignant par $\varphi_M(b)$ l'homothétie dans $M$ définie par un $b\in B$ : $g\varphi_M(b)g^{-1}=\varphi_M(gb)$ d'où $\varphi_M(\operatorname{Tr}_G b)=\operatorname{Tr}_G\varphi_M(b)$, d'où le

% B05 | first Corollaire and bracketed qualification
\uline{Corollaire} Si $\lambda\in\operatorname{Tr}_G B$, l'homothétie $\lambda$ dans $\widehat H(G,M)$ est \uline{nulle} \inserted{[et la condition est \unsure{manifeste\textsuperscript{t}|manifestement} nécessaire, en prenant $M=\underline B$)} qui implique $a'\Longrightarrow b$. $b\Longrightarrow b(q)$ trivial. Or $b(q)\Longrightarrow b(q+1)$

% B06 | immediately before sequence
car on a

% B07 | at fold, exact sequence and assignment below
\[
\begin{aligned}
0\to M&\to M\otimes_Z Z(G)\to N\to 0\\[-2pt]
m&\to\textstyle\sum m\otimes g
\end{aligned}
\]

% B08 | four lines under sequence
où on met sur le terme milieu la \deleted{\unsure{$B$-}}structure de $B$-module définie par $\lambda(m\otimes g)=\lambda m\otimes g$, et la structure de $G$-module définie par $g'(m\otimes g)=g'm\otimes g'g$, alors $g'.(\lambda[m\otimes g])=(g'.\lambda)(g'.[m\otimes g])$ car les deux mbres sont \unsure{respt} $g'.(\lambda m\otimes g)=g'(\lambda m)\otimes g'g$ et $(g'\lambda)(g'm\otimes g'g)=(g'\lambda)(g'm)\otimes g'\unsure{g|m}$, c'est la $=$ chose car $g'.(\lambda m)=(g'.\lambda)(g'.m)$.

% B09 | dimension shift, down to second corollary
Alors on a, puisque $\widehat H(G,M\otimes_Z Z(G))=0$, $\widehat H^{q+1}(G,M)=\widehat H^q(G,N)=0$. On prouve de façon analogue $b(q)\Longrightarrow b(q-1)$, d'où $b(q)\Longrightarrow b(0)$. Mais $b(0)\Longrightarrow a$ trivialement.

% B10 | second Corollaire through bottom
\uline{Corollaire} Supposons $B$ \inserted{commutatif} intègre et intégralement clos, \inserted{$A$ local.} Les conditions de la proposition sont vérifiées si l'indice de ramification $e$ de $B$ sur $A$ ($={}$ ordre du groupe d'inertie) \deleted{\unsure{soit premier à la|est premier à la} \inserted{\unsure{divisible}} \unsure{premier à}} \deleted{la caractéristique \unsure{(supposée|résiduelle} $\ne0$) \unsure{par} revient à écrire un $u\in B$ tel que $\unsure{eu=1}$}. est inversible dans $A$, i.e. \unsure{non nul|une unité} dans $A/\unsure{\mathfrak m(A)}$ \deleted{\illegible{1}}

% archive sheet number: 41

\scanpage{43}

% B01 | top, opening two lines
Avec les conditions de la proposition générale, soit $G'\subset G$, et soit $A'=B^{G'}$, $A\subset A'\subset B$. Conditions équivalentes

% B02 | upper, relative trace
\[\begin{aligned}
a)\quad &\operatorname{Tr}_{G/G'} A'=A \qquad(\operatorname{Tr}_{G/G'}=\text{transfert})\\
a')\quad &\operatorname{Tr}_{G/G'} A'\ni1
\end{aligned}\]

% B03 | upper, two maps and two empty labels
\[\begin{aligned}
b)\quad &\text{Pour tout $(G,B)$-module $M$,}\quad \widehat H(G',M)^{\deleted{G}}\xrightarrow[\deleted{\unsure{t|t(G,G')}}]{t(G,G')}\widehat H(G,M)^{\deleted{\unsure{G|1}}}\ \text{est \deleted{injectif} \inserted{bijectif}}\\
b')\quad &\deleted{\unsure{\widehat H(G|\widehat H(G'}}\widehat H(G',M)\longrightarrow\widehat H(G,M)\quad\text{\deleted{bijectif} \inserted{surjectif}}\\
b(q))\quad &\\ b'(q)\quad &
\end{aligned}\]

% B04 | middle, under qualifying conditions
$a)\Leftrightarrow a'$ trivialement. $a')\Rightarrow b)$ car on aura $t(G,G')\varphi(b)=\varphi(t(G,G')b)$ pour $b\in A'$, d'où si $t(G,G')b=1$, $t(G,G')\varphi(b)=\underline{1}$, d'où $b)$ d'après un lemme général.

% B05 | middle right, two-tier implications
\diagram{\begin{tikzcd}[ampersand replacement=\&]
 b \arrow[r,Rightarrow] \arrow[d,Rightarrow] \& b' \arrow[d,Rightarrow]\\
 b(q) \arrow[r,Rightarrow] \& b'(q)
\end{tikzcd}}

% B06 | left, oblique beside middle of proof
% slanted
\margin{\unsure{On peut l'appliquer} \unsure{encore|en effet} en prenant $G'=G_i$, mais aussi pour $G'$ un sous-groupe de Sylow $H$ pour $p$, où $p$ est la \uline{caractéristique} du $\overline k$. Les deux reviennent à étudier $(\unsure{O|\mathcal O}'{}^{H}/\operatorname{Tr}_H O')^{G}$, donc d'abord à déterminer un $O'{}^{H}/\operatorname{Tr}_H O'$ pour une $p^{\unsure{h|n}}$-extension.}

% B07 | middle, dimension shift and degree zero
De plus on voit toujours de la $=$ façon que $b'(q)\Rightarrow b'(q+1)$ et $b'(q)\Rightarrow b'(q-1)$, donc $b'(q)\Leftrightarrow b'(0)$. \unsure{Appliquons|Appliquant} $b'(0)$ à $M=B$, on trouve que $\deleted{\unsure{\widehat H|A}} A/\operatorname{Tr}_G B\longleftarrow(A'/\operatorname{Tr}_{G'} B)^{\deleted{\unsure{G}}}$ est surjectif, i.e. que l'on a $1\in\operatorname{Tr}_G B+\operatorname{Tr}_{G/G'}A'=\operatorname{Tr}_{G/G'} A'$.

% B08 | lower, underlined heading and lemma
\uline{Reste à prouver} Soient $G'\subset G$ ($G$ fini) soit $M$ un $G$-module, tel que $1\in\operatorname{Hom}_{\mathbb Z}(M,M)$ soit du type $1=\operatorname{Tr}_{G/G'}\varphi$ \deleted{\illegible{1}} ($\varphi\in\operatorname{Hom}_{G'}(M,M)$) alors

% B09 | lower middle, opposing arrows
\[\begin{aligned}
\operatorname{Tr}_{G/G'}\qquad &\widehat H\deleted{\unsure{G|M}}(G',M)\longrightarrow\widehat H(G,M)\quad\text{est surjective}\\
\operatorname{Tr}'_{G'/G}\qquad &\widehat H(G',M)\longleftarrow\widehat H(G,M)\quad\text{est injection,}
\end{aligned}\]

% B10 | lower, stable elements
et l'image de ce\deleted{\illegible{1}} dernier homo est l'ens des éléments de $\widehat H(G',M)$ qui sont ``stables'', \unsure{d.té|dé.té} $\widehat H(G',M)^G$.

% B11 | lower left, enclosure linked to transfer
\margin{Le coimage du premier est $\widehat H(G',M)_G$, i.e. le quotient de $\widehat H(G',M)$ par ses \unsure{sous|syst} \deleted{\unsure{non stables}} \unsure{puissances|parties}}

% B12 | foot, Sylow remark
(Ceci contient le th. sur les groupes de Sylow, en plaçant dans la catégorie des groupes abéliens \deleted{\unsure{de torsion}} \inserted{mod} les groupes de torsion où $p$ est est inversible.

% B13 | very foot, crosses below final rule
\unsure{Avec compatibilité|Avec égalité} entre les \unsure{différents|coefficients} $\deleted{\unsure{\operatorname{Tr}_{G/G'}\,g}}\operatorname{Tr}_{G'/G'\cap G':g}\,\{\,\deleted{\unsure{g}}g\,\operatorname{Tr}_{G':g/G'\cap G':g}\,\}$.

% archive sheet number: 42

\scanpage{44}

% B01 | head, item 1
1) Soit $G$ groupe fini d'automorphismes d'un anneau $1/2$ local $O$. $G$ est transitif sur les idéaux maximaux $\underline m$, soit $G_d$ le stabilisateur de $\underline m_0$, et $G_i\subset G_d$ l'ens des élt de $G_d$ qui définissent l'automorphisme identique de $O/\underline m$. Sous ces conditions

% B02 | upper, part a
a) Il existe $x\in O^{G_i}$ tel que $T_{G/G_i}\,x=1$. Donc pour \unsure{tout|$\forall$} $G$-$O$-module $M$, on a

% B03 | upper, stable elements formula
\[\widehat H(G,M)=\widehat H(G_i,M)^{G/G_i}\]

% B04 | next three lines
Donc pour qu'on ait $1\in T_G O$, i.e. $\widehat H(G,M)=0$ ql soit $M$, il f. et s. que $[G_i:e]$ ne soit pas multiple de la caractéristique $p$ de $O/\underline m$.

% B05 | middle, part b
b) Si $G_i=\{e\}$, on a $H^1(G,O^*)=\{e\}$.

% B06 | middle, question beside vertical rule
Plus précisément, quel que soit $G_i$, l'hom naturel $H^1(G,O^*)\longrightarrow H^1(G_i,(\underline O\deleted{\unsure{G|G_i}}/m)^*)$ est bijectif (?). Qu'en est-il dans le cas \deleted{\unsure{G}} des $\operatorname{Gl}(n,\underline O)$ au lieu de $O^*$.

% B07 | left of restriction question
\margin{?}

% B08 | low left, diagonal
% slanted
\margin{dans le cas où \unsure{$O$|$A$} est de dim $1$, car alors $\widehat H(G,M)$!}

% B09 | lower third, after blank
Pour $a$, noter

% B10 | lower third, overlines and sum
\[\begin{aligned}
\overline{\widehat H(G,M)}&=\widehat H(G,\overline M)=\widehat H(G,\bigoplus_{s\in G/G_d}s.\overline M_{\underline m_0})=\widehat H(G_d,\overline M_{\underline m_0})\\
&=\overline{H(G_d,M_{\underline m_0})}.
\end{aligned}\]

% B11 | beneath first equality, brace
\inserted{justification par Zariski; il y a donc hyp. noethérienne \unsure{finie??|fausse??|faire??}}

% B12 | above and left of diagram
Considérons alors le diagramme :

% B13 | foot, left and right vertical maps
\diagram{\begin{tikzcd}[ampersand replacement=\&,column sep=small]
\widehat H(G,M) \arrow[r,"\sim"] \arrow[d,"\wr"'] \& \widehat H(G_i,M)^{G/G_i}_{\deleted{\unsure{G|G_i}}} \arrow[r,"\deleted{\unsure{g|\mathrm{res}}}","\alpha"'] \& \widehat H(G_i,M_{\underline m_0})^{G_d/G_i}\\
\widehat H(G_d,M)^{G/G_d} \arrow[rr,"\beta"] \& \& \widehat H(G_d,M_{\underline m_0})\simeq\widehat H(G_d,M)_{O_d-m\cap O_d}\arrow[u,"\wr"]
\end{tikzcd}}

% B14 | pencil, by enclosed right-hand column
\margin{on passe aux \unsure{complétés|quotients}!}

% B15 | bottom right, upright rule
\margin{$\alpha$ et $\beta$ sont injectifs, à images denses}

% archive sheet number: 43

\scanpage{45}

% B01 | upper, first exact row, captions beneath objects
\diagram{\begin{tikzcd}[ampersand replacement=\&,column sep=small]
0\arrow[r] \& G_k^0\arrow[r] \& \underset{\substack{\text{extension}\\\text{abélienne}\\\text{maximale}}}{G_k}\arrow[r] \& \underset{\substack{\text{extension}\\\text{non ramifiée}\\\text{maximale}}}{\unsure{\mathfrak P|P}_k}\arrow[r] \& 0
\end{tikzcd}}

% B02 | upper, second exact row, labels in ink and pencil
\diagram{\begin{tikzcd}[ampersand replacement=\&,column sep=small]
0 \arrow[r] \& Z_2^{r_1}\arrow[r] \& G_k^0\arrow[r] \& \underbrace{U_k^0}_{\substack{\text{idèles}\\\text{unités}\\\text{finis}}}/\underbrace{\operatorname{Im} E}_{\substack{\text{idèles}\\\text{unités}\\\text{finis}\\\text{principaux}}}\arrow[r] \&0
\end{tikzcd}}

% B03 | pencil below group in second row
\inserted{\unsure{on en fait quoi ???|on en fait vraiment ???|on en fait tant ???}}

% B04 | centre, one large oblique cross over N.B. and following sequences
\deleted{N.B. $G_k/Z_2^{r_1}$ \deleted{qui} est un groupe dont la définition ne fait plus intervenir d'idèle à l'infini) est le groupe de Galois de l'extension abélienne non ramifiée à l'infini maximale (i.e. qui est \unsure{complexe|couverte} au dessus de tte place réelle de $k$) -- tout de façon générale, $S$ un ensemble \inserted{fini} de places finies ou infinies réelles, alors le groupe de Galois top. de l'extension \inserted{abélienne} non ramifiée \unsure{ds|en} $\complement S$ maximale, est $G_k^S$, où}

% B05 | lower, cancelled exact rows, same cross
\diagram{\begin{tikzcd}[ampersand replacement=\&,column sep=small]
\deleted{0}\arrow[r] \& \deleted{G_k^{\unsure{0S|0\underline S|0,S}}}\arrow[r] \& \deleted{G_k^S}\arrow[r] \& \deleted{\mathfrak P_k}\arrow[r] \& \deleted{0}\\
\deleted{0}\arrow[r] \& \deleted{Z_2^{S^\infty}/\bigcap_k\operatorname{Im}((E^{S^\infty})^k)}\arrow[r] \& \deleted{G_k^{0\underline S}}\arrow[r] \& \deleted{U_k^{S^0}/\operatorname{Im}E}\arrow[r] \& \deleted{0}
\end{tikzcd}}

% B06 | underbrace below denominator in final exact row
\inserted{ceci est-il identique à $Z_2^{S^\infty}$ ???}

% B07 | bottom, underlined lemma heading
Ceci résulte du \uline{Lemme} Pour qu'une \deleted{\unsure{extension|abélienne}}\inserted{extension} $L$ d'un corps local \deleted{\unsure{soit}} quel $K$ soit non ramifiée il faut et il suffit que $N_{L/K}U_L=U_K$ // ou $N_{L/K}L^*\supset U_K$

% archive sheet number: 44

\scanpage{47}

% B01 | head, underlined proposition
\uline{Proposition} Soit $k$ un corps, $n$ un entier. Conditions équivalentes

% B02 | upper, clauses a and b
a) $H^i(k,G)=0$ pour $i>n$ et tout groupe alg. comm. $G$ déf$/k$

% B03 | upper, clause b
b) $H^i(k,\Gamma\deleted{\unsure{G|0}})=0$ \inserted{\uline{fini} $\Gamma$ déf$/k$}

% B04 | upper, clause b prime, qualification above
b') $H^i(k,\Gamma)=0$ pour $i>n$ et tout \deleted{groupe espace vectoriel} \inserted{groupe \inserted{fini} $\Gamma$ qui est un} esp. vectoriel sur un corps $Z/\ell$, dans lequel $\mathfrak g_k$ opère (avec noyau ouvert) irréductiblement

% B05 | upper left, oblique at b and b prime
% slanted
\margin{$b,b'$ sont tjrs vérifiés $(\unsure{n\geq1|n>1})$ si $\Gamma$ est un $p$-groupe \deleted{\illegible{1}}}

% B06 | upper middle, clause c, equals with no right operand
c) $H^i(k,G)=$ pour $i>n$ et tout groupe alg. comm. $G$ déf$/k$ \deleted{connexe et affine} devenant isomorphe à un groupe \deleted{\unsure{tore|torus}} $(G_m)^n$ sur la clôture galoisienne de $k$

% B07 | upper middle, clause c prime, left blank and insertion
c') $H^i(k,G)$ \deleted{\unsure{pr|que}} pour $i>n$, chaque fois que $G$ est le groupe des éléments inversibles d'une algèbre \inserted{comm.} \deleted{semi-simple} \inserted{séparable} déf$/k$

% B08 | just below fold
\uline{Démonstration}

% B09 | middle, four branches of implication
\diagram{\begin{tikzcd}[ampersand replacement=\&,row sep=small]
\&a\arrow[dl,Rightarrow]\arrow[dr,Rightarrow]\&\\
 b\arrow[d,Rightarrow]\&\&c\arrow[d,phantom,"\deleted{\Downarrow}" description]\\
 b'\&\&\unsure{a'|c'}
\end{tikzcd}}

% B10 | middle, proof opening
est trivial. Nous prouverons $c\Rightarrow b'$ et $b'\Rightarrow\deleted{b} a$ ce qui achèvera la démonstration

% B11 | below, first reduction in brackets
$b'\Rightarrow b)$ \deleted{En effet $b'\Rightarrow b$} comme on voit par un dévissage immédiat [\deleted{supposant} \inserted{se ramenant} d'abord à $\Gamma$ \deleted{d'ordre} \inserted{annulé par} \inserted{un} $\ell^k$, puis en filtrant par $\Gamma\supset\ell\Gamma\supset\ell^2\Gamma\supset\ldots\supset\ell^k\Gamma=0$, au cas où $\Gamma$ est annulé par $\ell$ i.e. est un vectoriel sur $Z/\ell$, puis enfin au cas irréductible].

% B12 | bottom, fresh implication and struck start
$b\Rightarrow a$ \deleted{En effet considérons le plus grand sous-groupe linéaire $G_1$ de $G$, (qui est défini sur $k$). Alors \illegible{2}} Supposons d'abord que le corps \deleted{soit \`a} $k$ soit \uline{parfait} Soit alors $G_1$ le plus grand sous-groupe linéaire de $G$, on sait qu'il est \uline{défini sur $k$}, et on a la suite exacte

% B13 | foot, exact sequence
\[0\longrightarrow G_1\longrightarrow G\deleted{\unsure{G_1}}\longrightarrow G/G_1\longrightarrow0\qquad(\text{déf}/k)\]

% archive sheet number: 46

\scanpage{48}

% B01 | head, reduction and decomposition
Cela nous ramène au cas où $G$ est soit complet, soit linéaire. Dans le \deleted{premier} \inserted{deuxième} cas, \deleted{considérons pour tout entier $n$} $G=G_u\times G_{ss}$ (décomposition en partie semi-simple et unipotente) et cette décomposition est \uline{définie sur $k$}.

% B02 | left of underlined descent assertion
\margin{à vérifier}

% B03 | upper middle, cases and torsion
Donc on peut supposer que $G$ est soit unipotent, soit \deleted{semi-simple} formé d'éléments semi-simples, soit complet. Dans les deux derniers cas, soit $U$ l'ens. des éléments d'ordre fini de $G=G^{\bar k}$; alors le groupe $G^{\deleted{\unsure{\bar k|k}}}/U$ est un groupe sans torsion, de plus il est divisible; pour le voir, on note que $G_0$ (comp. connexe de $0$) est divisible (c'est le th. de Weil pour une variété abélienne, et \deleted{bien sûr} trivial pour le cas où $G\simeq\bar k^{*\,n}$) et \deleted{\unsure{donc|th}} par suite $G\simeq G_0+G/G_0$, \deleted{\unsure{si|et}} (abstraction faite des opérations de $\mathfrak g_k$) d'où résulte que $U$ contient le facteur $G/G_0$ et que $G/U$ s'identifie au \deleted{\illegible{1}} quotient de $G_0$ par $U\cap G_0$, donc est divisible aussi).

% B04 | left centre, note in two directions
% slanted
\margin{lemme : dégager cf plus \unsure{bas|loin} ici le fait de passer à la limite inductive est essentiel}

% B05 | below, consequence
Donc $\deleted{\unsure{H}}\underline{H}^i(\mathfrak g,G/U)=0$, \inserted{pour $i>0$} \deleted{tandis que} d'où

% B06 | lower quarter, comparison isomorphism
\[\underline{H}^i(\mathfrak g,G)\simeq H^i(\mathfrak g,U)\qquad\text{pour }i\geq1\]

% B07 | lower left, oblique beside comparison
% slanted
\margin{dégager un \uline{lemme}, [lorsque $G$ n'a pas de sous-groupe linéaire \unsure{connexe|commun} unipotent non nul]}

% B08 | lower, parentheses and transition
(compte tenu que $G\longrightarrow G/U$ est surjectif.) Or, pour \unsure{tout|$\forall$} $m$, soit $U(m)$ le sous-groupe de $U$ formé des éléments annulés par $m$. On a

% B09 | foot, direct limit
\[H^i(\mathfrak g,U)=\lim_{\underset m\longrightarrow} H^i(\mathfrak g,U(m))\]

% archive sheet number: 47

\scanpage{49}

% B01 | top, continuation of direct limit
et lorsque $i>n$, les \deleted{\unsure{de cette|derniers|deux}} $H^i(\mathfrak g,U(m))$ sont nuls d'après $b)$, \inserted{(car $U(m)$ est fini)} d'où résulte notre assertion.

% B02 | upper quarter, unipotent group
Si enfin $G$ est \deleted{le groupe} un groupe linéaire unipotent connexe, on distingue deux cas

% B03 | upper middle, characteristic zero
1°) Caractéristique nulle. Alors $G$ est $\simeq\bar k^n$ sur $\bar k$, donc aussi sur $k$ (en vertu de $H^1(k,\operatorname{Gl}(\unsure{n|m}))=0$). Et comme $H^i(k,G_a)=0$ pour $i>0$, l'assertion est valable.

% B04 | middle, nonzero characteristic
2°) Car. $p\ne0$. Alors on filtre $G$ par $G\supset pG\supset\ldots\supset p^\ell G=0$. On est ramené au cas où $pG=0$ (les groupes $p^iG$ étant \uline{définis sur $k$}).

% B05 | left, opposite pG = 0
\margin{$//$}

% B06 | middle, descent
Mais on sait qu'alors $G$ est $\simeq\bar k^{\unsure{m|n}}$ sur $\bar k$ Alors \deleted{\unsure{=|\simeq}} aussi sur $k$ (car $H^1(k,\operatorname{Gl}(\unsure{m|n}))=0$), et on finit comme dans 1°.

% B07 | last line, followed by blank half-page
Supposons que $k$ soit quelconque .....

% archive sheet number: 48

\scanpage{50}

% B01 | head, proof of implication and inserted hypothesis
$c\unsure{'|)}\Rightarrow b'$ Soit en effet $\Gamma$ un groupe comme dans $b')$, \inserted{Supposons d'abord $\ell\ne p$.} Soit $A=\bar k(\Gamma)$ l'algèbre de $\Gamma$ à coefficients dans la clôture algébrique $\bar k$ de $k$, \inserted{$\mathfrak g_k$ y opérant de façon évidente.} C'est une algèbre déf$/k$, commutative et \deleted{semi-simple} sép\uline{ara}ble (car $[\Gamma:e]$ est premier à $p$), donc en vertu de $c')$, on a $H^i(k,G)=0$ si $i>n$, $G$ désignant le groupe des éléments inversibles de $A$.

% B02 | upper middle, norm and its kernel
Soit $G_1$ le sous-groupe de $G$ formé des $x\in G$ tels que $N_{A/\bar k}x=1$ (où \deleted{dét. \unsure{$=$|est}} $N_{A/k}x=$ dét. de la translation à g. par $x$ dans $A$), c'est un sous-groupe \uline{défini sur $k$} (noyau de l'homomorphisme $N_{A/k}:G\longrightarrow\bar k^*$, qui est déf$/k$).

% B03 | middle, simple module and intersection
Comme $\Gamma$ est \uline{simple} comme $\mathfrak g$-module, \inserted{et que $\Gamma\subset G$} il en résulte que \deleted{\unsure{$N_{A/k}$}} l'on a $\Gamma\subset G_1$, ou $\Gamma\cap G_1=\deleted{\unsure{\{1\}|\{0\}}}\{e\}$, et on vérifie facilement que ce deuxième cas est exclu (distinguer le cas \deleted{\unsure{pair|$p$|premier}}\inserted{$\ell$} impair, et le cas $\ell=2$ $\deleted{\unsure{8|\}}}$, on note \deleted{\unsure{que}} dans ce dernier cas qu'en tous cas on a $N_{A/k}\xi=\pm1$ si $\xi\in\Gamma$).

% B04 | lower, order of the multiplication kernel
Donc $\Gamma\subset G_1$. Comme $\ell\Gamma\deleted{\unsure{G|=0}}=\{0\}$, il en résulte que $\Gamma$ est contenu dans le noyau de l'hom. $G_1\xrightarrow{\ell}G_1$. Ce noyau est d'ailleurs d'ordre $\ell^{n_1}$, où $n_1=n-1$ est la dim. de $G_1$, donc $n$ l'ordre de $G$.

% archive sheet number: 49

\scanpage{51}

% B01 | head, lemma and vertical rule
\uline{Lemme 1} Soit $k$ un corps algébriquement clos, $K$ une extension de type fini, de degré de transcendance $1$, de $k$. Alors $K$ est \deleted{\unsure{algèb|algéb}} quasi-algébriquement clos

% B02 | upper, first corollary
\uline{Corollaire 1} $\widehat H^i(K^{\unsure{*|2|'}},G_m)=0$ pour tout $i$

% B03 | upper, proof
En effet, $K$ et ses extensions galoisiennes $K'$ sont quasi-alg. clos, donc \deleted{\unsure{$H^i(K,G$}} la conclusion est bien connue.

% B04 | upper middle, second corollary
\uline{Corollaire 2} $\deleted{\widehat{\vphantom H}}H^i(K,\Gamma)=0$ \deleted{si $i>0$} \inserted{pour tt $i>1$ et} \deleted{pour} tout groupe abélien fini $\Gamma$ où $G_K$ opère trivialement.

% B05 | middle, roots of unity
Soit en effet $\widetilde K$ \deleted{la} \inserted{une} clôture galoisienne de $K$, [$U\subset\widetilde K^*$ le \deleted{\unsure{sou|un|sous-}} groupe des racines de l'unité dans $\widetilde K$. Alors $\widetilde K^*/U$ est uniquement divisible sauf pour $p$ (\unsure{l'exposant|le apparent} car.) d'où la suite exacte \uline{mod} $p$

% B06 | middle, comparison across the fold
\[\deleted{\unsure{\simeq|=}}\widehat H^i(K,\deleted{\unsure{G_m|\widetilde K}}U)\xrightarrow{\sim}\widehat H^i(K,G_m)\deleted{\unsure{\mathrm{est}|0}}\qquad\text{\uuline{mod $p$}}\]

% B07 | left of comparison, two slightly slanting lines
% slanted
\margin{a-t-on un vrai isomorphisme?}

% B08 | below fold, nullity and its revision
Or le \deleted{premier} deuxième mb est nul, donc le premier l'est, donc comme le premier \deleted{\unsure{sing|groupe}} est $p$-composante nulle, il est nul.] On \deleted{\unsure{prouve|pourra}} \deleted{\unsure{en vertu de|en vertu}} voit plus généralement que si $n$ est premier à la caractéristique, alors

% B09 | lower third, power map
$\widetilde K\xrightarrow{x\to x^n}\widetilde K^{\unsure{n|\times}}$ est un épimorphisme de \deleted{groupe de} noyau $\simeq\mathbb Z/n$, donc la suite exacte de cohomologie donne le résultat $H^i(K,\mathbb Z_n)=0$.

% B10 | foot, deleted heading and characteristic p
\deleted{\uline{Lemme 2} Soit} Reste à prouver $H^i(K,\mathbb Z/p)=0$ si $i>1$. Pour cela on considère la suite exacte $0\longrightarrow\mathbb Z/p\longrightarrow\widetilde K\xrightarrow{\wp}\widetilde K\unsure{)|'}\longrightarrow0$.

% archive sheet number: 50

